Cooperative aerial transport involves $N$ multirotor uncrewed aerial vehicles (UAVs) lifting a common payload via flexible cables. This platform class has been extensively studied and demonstrated in laboratory and structured outdoor experiments for disaster-response resupply, suspended construction, and heavy-payload scenarios~\cite{Michael2011, Fink2011, Bernard2011, Tagliabue2019}. This work focuses on the operational regime in which the platform's limiting behavior is determined by its fault response. Cable severance may result from abrasion, attachment-hardware fatigue, or obstacle contact; the tension previously transmitted by the severed cable redistributes across the remaining cables within one payload-pendulum period. If a controller cannot absorb this redistribution within that timeframe, the payload is dropped. In contrast to actuator or sensor faults in single-vehicle systems, cable severance constitutes a structural change in the plant: it is a discrete reduction of the constraint set rather than an additive perturbation. This event occurs without any supervisory signal and necessitates a corrective response on a timescale dictated by the payload pendulum, typically well under one second.

Three structural features of this plant—unilateral compliant coupling~\cite{Brogliato2016}, hybrid constraint-cardinality transitions at severance events~\cite{GoebelSanfeliceTeel2012}, and strict information locality (each vehicle observes only its own state, rope tension, and a local payload estimate)—collectively preclude the use of classical cooperative and slung-load architectures. Any admissible controller must provide post-fault response without detection, annunciation, or supervisory reconfiguration, and must support a closed-loop analysis that remains valid across the structural transitions.

This paper demonstrates that, for this plant class, fault tolerance can be achieved through a single feed-forward identity in the per-drone controller, rather than through detection and subsequent reconfiguration: the altitude thrust feed-forward of drone $i$ is set equal to its own measured rope tension, $T_i^{\mathrm{ff}}(t) = T_i(t)$ (formally \eqref{eq:Tff} in \cref{sec:proposed}).

Here, $T_i^{\mathrm{ff}}$ denotes the thrust feed-forward added to drone $i$'s commanded acceleration, and $T_i$ represents its measured rope tension. The underlying rationale is based on timescale: when a peer's cable severs, the load redistributes across the remaining cables within one payload-pendulum period (approximately $2$~s), during which the surviving drone's measured tension increases. The identity \eqref{eq:Tff} translates this increase at the next control tick into a proportional thrust increment, thereby absorbing the new load at the plant's native response rate, without requiring detection or communication. This mechanism is embedded within a conventional cascade comprising an outer-loop proportional-derivative (PD) controller with anti-swing slot shift, a single-step convex quadratic program (QP) on the tilt-and-thrust envelope, and a geometric attitude loop. The identity \eqref{eq:Tff} is the sole nonstandard component.

This work establishes the consequences of \eqref{eq:Tff} along three axes. The first contribution is a taut-cable reduction theorem (C1, \cref{sec:reduction-and-reference}): within a slack-excursion-bounded domain, defined by a pendulum-to-rope timescale separation ratio $\delta$ and a bound $\eta$ on the fraction of time any cable is slack, an explicit first-order error budget is derived to justify replacing the rope's distributed spring-damper dynamics with a lumped effective stiffness. The domain's admissibility gates (slack-run duration $\le \SI{40}{\milli\second}$, duty cycle $\le 2.5\%$, QP active-set-transition fraction $\le 1\%$) are satisfied by all six demonstrated missions with margin; the worst observed values are $\SI{35.8}{\milli\second}$, $1.66\%$, and $0.10\%$, respectively (\cref{tab:domain-audit}). The second contribution is a baseline hybrid practical input-to-state stability (ISS) theorem (C2, \cref{sec:stability}): on the domain of C1, under a single-active-regime doctrine for the QP active set and a spaced-fault dwell-time hypothesis requiring faults to be separated by at least one payload-pendulum period, the closed-loop tracking error admits an explicit recovery envelope with a Lyapunov decay rate $\lambda$ and a per-fault-cycle Lyapunov contraction coefficient $\rho \in (0,1)$, both in closed form. The strict inequality $\rho < 1$ is achieved specifically by the identity \eqref{eq:Tff}, which certifies Lyapunov contraction across each fault cycle rather than mere non-divergence within a pendulum period. The dwell-time requirement $\tau_d \ge \tau_{\mathrm{pend}}$ is structural rather than conservative: it is the minimum time over which the measured-tension feed-forward term $T_i^{\mathrm{ff}} = T_i$ can restore the Lyapunov function below its pre-fault level with $\rho < 1$ (Lemma~\ref{lem:dwell-decay}). The third contribution is a closed-form stability bound for the $L_1$-adaptive altitude augmentation (C3, \cref{sec:extensions}): for the discrete-time $L_1$ outer loop on the altitude channel, a certifiable gain-selection condition is derived in closed form in terms of the control-tick period and the altitude-channel entry of the outer-loop Lyapunov matrix, thereby eliminating the dominant hyper-parameter of the $L_1$ architecture from empirical tuning.

The proposed controller is validated in high-fidelity Drake multibody simulation through a pre-declared six-variant capability demonstration and four targeted stress campaigns on a five-drone, \SI{10}{\kilo\gram}-payload, \SI{4}{\metre\per\second} Dryden-wind scenario under both single and compound cable-severance faults. The feed-forward ablation campaign isolates the identity \eqref{eq:Tff} as the causal mechanism: disabling this identity increases cruise RMS tracking error by $34$--$39\,\%$ and peak post-fault payload altitude sag by a factor of $3.6$--$4.0$. Two null findings are reported alongside the positive results: in these regimes, a composable extension layer remains structurally inactive because the baseline already satisfies the relevant constraint. These null results are included in the pre-declared protocol. The architecture admits a sequence of $F \le N - 2$ unannounced severances, subject to the actuator-margin envelope (\cref{sec:fault-and-controller}); the demonstrated configuration at $N=5$, $m_L=\SI{10}{\kilo\gram}$ admits $F \le 2$.

\paragraph{Architectural properties as affirmative results.}
The contribution is properly read not as a feature combination on a five-drone case study but as a category claim about how to handle structural transitions whose information arrives through actuator-side measurements on the plant's native timescale. Four architectural properties of the proposed law are stated as the consequences of that thesis. \emph{Locality:} the per-drone admissible information set $\mathcal{I}_i$ \eqref{eq:info-set} contains no peer state, no peer tension, no fault flag, and no shared scheduler; each drone executes its controller in isolation. \emph{Passivity (no detector, no reconfiguration):} the same identity \eqref{eq:Tff} that operates pre-fault produces the post-fault response, with no observer dynamics, no detection threshold, no classification logic, and no reassignment rule entering the closed loop. \emph{Robustness to multi-agent failure:} because the per-drone closed loop is identical pre- and post-fault and contains no peer-coupled term, the system tolerates a sequence of $F$ unannounced cable severances --- up to the geometric/actuator bound $F \le N - 2$ with $m_L g/(N{-}F) \le \kappa_{\mathrm{act}} f_{\max}$ \eqref{eq:actuator-envelope} --- by induction on the fault counter, with the recovery envelope of \cref{thm:hybrid-stability} carrying unchanged at each event; the architectural cost of an additional failure is one more bounded jump, $\chi(\Delta_f)$, which the FF identity controls and the closed-loop spectrum exponentially contracts within one inter-fault dwell. \emph{Team-size independence:} the per-drone law and the lemma stack of \cref{sec:stability} carry no $N$ structurally; team-size dependence is concentrated in two scalar quantities --- the per-drone hover load $m_L g/N$ and the post-fault redistribution $m_L g(1/(N{-}F) - 1/N)$ --- which together define the actuator-margin envelope above within which the certificate is invariant in $N$. The two null findings (\cref{rem:mpc-null,rem:reshape-null}) are read in this framing as evidence \emph{that the conventional active fault-tolerance (FT) machinery is not what is doing the work}: the model-predictive-control (MPC) and reshape layers, both passively present, are structurally inactive on the demonstrated trajectory family precisely because the underlying feed-forward identity already satisfies the relevant constraints.


The remainder of the paper is organized as follows. Section~\ref{sec:Problem_Statement} formalizes the plant, the fault model, the admissible controller class, the reference trajectory, the taut-cable reduction theorem, and the exact problem addressed. Section~\ref{sec:proposed} presents the proposed decentralized controller and its three-layer cascade. Section~\ref{sec:stability} establishes the main hybrid practical-ISS theorem through five supporting lemmas and characterizes its robustness. Section~\ref{sec:extensions} presents parameter adaptation via $L_1$-augmentation of the altitude channel, receding-horizon MPC with penalized tension-ceiling constraints, and the post-fault formation-reshape supervisor, and demonstrates that each preserves the baseline theorem's recovery envelope on its admissibility domain.

\subsection{Related Work}
\label{sec:related-work}

Existing work relevant to this problem can be grouped into four categories: centralized and distributed cooperative transport, geometric tracking and slung-load control, $L_1$ adaptive control, and fault-tolerant control. The following review maps these four categories and identifies the specific gap this paper addresses.

\paragraph{Centralized and distributed cooperative transport.} Force-allocation architectures originating with Michael, Fink, and Kumar~\cite{Michael2011} solve a centralized quadratic program over cable tensions to satisfy wrench-closure and actuator-limit constraints; subsequent refinements improve robustness and payload-mass estimation~\cite{Fink2011, Tagliabue2019} but preserve synchronous global state. Distributed predictive control~\cite{WangOng2017} replaces per-tick global communication with iterative neighbor exchanges, but message complexity still scales with the horizon and the coupling rank. More recent work extends nonlinear MPC to six-degree-of-freedom rigid-body payload transport~\cite{LiLoianno2023}, adds control-barrier inter-vehicle safety~\cite{PallarLi2025}, and handles parametric payload uncertainty across directed communication topologies~\cite{SuBhowmick2023}. At agile multirotor tick rates, the per-tick synchronous global-state requirement common to all of these formulations is a structural design constraint that prevents their application at tick rates where cable-severance events can occur within a single outer-loop horizon. None of these architectures accounts for cable-severance events in its stability argument; each assumes a fixed constraint set over the control horizon.

\paragraph{Geometric tracking and slung-load control.} The foundational single-quadrotor geometric tracking controller~\cite{LeeLeokMcClamroch2010}, the flatness-based slung-load framework~\cite{Sreenath2013}, and the $\mathrm{SE}(3)$ geometric control for a rigid-body suspended payload~\cite{Lee2018} all sit under a fixed-constraint-set assumption. Multi-vehicle extensions~\cite{LeeSreenathKumar2013, SharmaSundaram2023} reintroduce rigid-formation kinematics or global-state exchange, both of which collapse the decentralization that motivates the multi-vehicle architecture.

\paragraph{$L_1$ adaptive control.} $L_1$ adaptive control~\cite{CaoHovakimyan2008, Hovakimyan2010book} delivers matched-disturbance rejection with transient-performance bounds, recently extended to the full-$\mathrm{SE}(3)$ dynamics of a single quadrotor~\cite{WuCheng2025}; the extension to cooperative architectures under cable-fault events and strict locality remains open.

The fault-tolerant control literature divides into active architectures that reconfigure after diagnosis~\cite{Blanke2015, ZhangJiang2008} and passive architectures that absorb faults through robustness~\cite{Patton1997}. Within cooperative aerial transport, two prior strands are directly relevant. First, fault-tolerant cooperative transportation under suspension-cable failures has been addressed through observer-based disturbance estimation with dynamic load reassignment~\cite{LiuSuspensionFailure2021}: each surviving vehicle estimates a tension-related disturbance, and the load allocation is redistributed through a reassignment law. Second, communication-free or low-communication cooperative transport is achieved by passive force control~\cite{GassnerPassiveForce2017} and leader–follower architectures with force estimation~\cite{NoCommTransport2021}, which avoid inter-UAV message exchange but presuppose a fixed cable-cardinality regime. These strands demonstrate that neither cable-severance response nor strict locality is independently novel in this field. The contribution of the present work is best read not as a feature subtraction relative to existing methods but as a different commitment about \emph{how the structural-transition announcement is processed}. Against the observer-based reassignment line~\cite{LiuSuspensionFailure2021}, the difference is not the absence of an observer but the absence of \emph{the announcement-detection step entirely}: the proposed architecture treats the actuator-side tension measurement as the announcement, eliminating the detection-threshold/confirmation-window timeline on which active reconfiguration depends. Against the passive force-control and leader--follower lines~\cite{GassnerPassiveForce2017, NoCommTransport2021}, the difference is the explicit hybrid treatment of the constraint-cardinality transition: the fixed-cardinality presupposition is replaced by an analysis that admits unannounced severances within the recovery envelope. Against the centralized and synchronous-global-state cooperative-transport literature~\cite{Michael2011, WangOng2017, LiLoianno2023, PallarLi2025, SuBhowmick2023}, the difference is that the announcement is consumed locally rather than aggregated globally --- the per-tick global-state synchronization that those formulations require is precisely the mechanism the proposed architecture avoids. The architectural category occupied by this paper, then, is decentralized cooperative manipulation in which discrete structural transitions self-announce through actuator-side measurements; the cable-severance instance is the certification vehicle, not the contribution.